\begin{filecontents*}[overwrite]{\jobname.ist}
headings_flag 1
numhead_positive "数字"
symhead_positive "符号"
heading_prefix "{\\bfseries\\hfil "
heading_suffix "\\hfil}\\nopagebreak\\par\\vskip 5pt "
\end{filecontents*}


\usepackage[a4paper,margin=28mm,headheight=15pt]{geometry}
\usepackage{amsmath,amssymb,mathtools,amsthm,bm}
\usepackage{graphicx}
\usepackage{tikz}
\usepackage{booktabs,longtable,array}
\usepackage{multicol}
\usepackage{enumitem}
\usepackage{xcolor}
\usepackage{fancyhdr}
\usepackage{imakeidx}
\makeindex[title={译本术语索引},intoc,noautomatic]
\usepackage{hyperref}

\renewcommand{\contentsname}{目录}
\renewcommand{\figurename}{图}
\renewcommand{\tablename}{表}
\ctexset{
  chapter = {
    name = {第,章},
    number = \arabic{chapter}
  }
}

\definecolor{jhcstblue}{HTML}{1F4E79}
\hypersetup{
  unicode=true,
  colorlinks=true,
  linkcolor=jhcstblue,
  citecolor=jhcstblue,
  urlcolor=jhcstblue,
  pdftitle={J-全纯曲线与辛拓扑（第二版）GLM \& GPT 译},
  pdfauthor={Dusa McDuff; Dietmar Salamon}
}

\setcounter{secnumdepth}{3}
\setcounter{tocdepth}{2}
\setlength{\parindent}{2em}
\setlength{\parskip}{0pt}
\emergencystretch=2em

\pagestyle{fancy}
\fancyhf{}
\fancyhead[LE]{\small\thepage\quad\nouppercase{\leftmark}}
\fancyhead[RO]{\small\nouppercase{\rightmark}\quad\thepage}
\renewcommand{\headrulewidth}{0.4pt}

\newtheorem{theorem}{定理}[section]
\newtheorem{proposition}[theorem]{命题}
\newtheorem{lemma}[theorem]{引理}
\newtheorem{corollary}[theorem]{推论}
\newtheorem{conjecture}[theorem]{猜想}
\newtheorem{question}[theorem]{问题}
\newtheorem*{namedtheorem}{定理}
\newtheorem*{namedproposition}{命题}
\theoremstyle{definition}
\newtheorem{definition}[theorem]{定义}
\newtheorem{example}[theorem]{例}
\newtheorem{exercise}[theorem]{习题}
\newtheorem*{nameddefinition}{定义}
\newtheorem*{namedexample}{例}
\theoremstyle{remark}
\newtheorem{remark}[theorem]{注}
\renewcommand{\proofname}{证明}

\numberwithin{equation}{section}

\DeclareMathOperator{\End}{End}
\DeclareMathOperator{\GL}{GL}
\DeclareMathOperator{\Ham}{Ham}
\DeclareMathOperator{\Hom}{Hom}
\DeclareMathOperator{\Vect}{Vect}

\newcommand{\zhindex}[2]{\index{#1@#2}}
\newcommand{\personfirst}[3]{#2（#3）\zhindex{#1}{#2（#3）}\ignorespaces}
\newcommand{\personlater}[3]{#2\zhindex{#1}{#2（#3）}\ignorespaces}
\newcommand{\srcite}[1]{\textup{[#1]}}

\newcommand{\sourcefootnote}[1]{\footnote{译者注：#1}\ignorespaces}

\providecommand{\gloss}[2]{\par\raggedright\hangindent1.5em\hangafter1\noindent #1\ \emph{#2}}
\providecommand{\idxentry}[2]{\par\raggedright\hangindent1.5em\hangafter1\noindent #1，#2}

\usepackage{multicol}
\newcommand{\BookGlossEntry}[2]{%
  \par\noindent\hangindent=1em\hangafter=1 #1\quad #2\par}
\newcommand{\BookSourceIndexEntry}[2]{%
  \par\begingroup\leftskip=#1em\relax\hangindent=1em\relax
  \noindent #2\par\endgroup}
\makeatletter
\newif\ifBookIndexing
\BookIndexingfalse
\newcommand{\BookBoldPage}[1]{\textbf{\hyperpage{#1}}}
\DeclareRobustCommand{\BookRecordIndex}[1]{%
  \ifBookIndexing
    \expandafter\index\expandafter{\csname BookIndex@#1\endcsname|hyperpage}%
  \fi}
\DeclareRobustCommand{\BookIndexDefinition}[1]{%
  \ifBookIndexing
    \expandafter\index\expandafter{\csname BookIndex@#1\endcsname|BookBoldPage}%
  \fi}
\DeclareRobustCommand{\BookTerm}[2]{%
  \BookRecordIndex{#1}%
  \begingroup\let\BookRecordIndex\@gobble #2\endgroup}
\AddToHook{cmd/@outputpage/before}{\let\BookRecordIndex\@gobble}
\pdfstringdefDisableCommands{%
  \def\BookTerm#1#2{#2}%
  \let\BookRecordIndex\@gobble
  \let\BookIndexDefinition\@gobble}
\expandafter\def\csname BookIndex@T0001\endcsname{omegajixianji@$\omega$-极限集（omega-limit set）}
\expandafter\def\csname BookIndex@T0002\endcsname{jbukefenjie@$J$-不可分解（J-indecomposable）}
\expandafter\def\csname BookIndex@T0003\endcsname{jquanchunquxian@$J$-全纯曲线（J-holomorphic curve）}
\expandafter\def\csname BookIndex@T0004\endcsname{jhanshu@$J$-函数（J-function）}
\expandafter\def\csname BookIndex@T0005\endcsname{kbubian@$k$-不变（k-invariant）}
\expandafter\def\csname BookIndex@T0006\endcsname{kdengbian@$k$-等变（k-equivariant）}
\expandafter\def\csname BookIndex@T0007\endcsname{mfenli@$m$-分离（m-separated）}
\expandafter\def\csname BookIndex@T0008\endcsname{nfenge@$n$-分割（n-splitting）}
\expandafter\def\csname BookIndex@T0009\endcsname{wfenge@$w$-分割（w-splitting）}
\expandafter\def\csname BookIndex@T0010\endcsname{amoxinglianluo@A 模型联络（A-model connection）}
\expandafter\def\csname BookIndex@T0011\endcsname{bradlowoudui@Bradlow 偶对（Bradlow pair）}
\expandafter\def\csname BookIndex@T0012\endcsname{calderonzygmundbudengshi@Calderón--Zygmund 不等式（Calderon--Zygmund inequality）}
\expandafter\def\csname BookIndex@T0013\endcsname{carlemanxiangsiyuanli@Carleman 相似原理（Carleman similarity principle）}
\expandafter\def\csname BookIndex@T0014\endcsname{donaldsonliangzifanchou@Donaldson 量子范畴（Donaldson quantum category）}
\expandafter\def\csname BookIndex@T0015\endcsname{dubrovinlianluo@Dubrovin 联络（Dubrovin connection）}
\expandafter\def\csname BookIndex@T0016\endcsname{eulerchang@Euler 场（Euler field）}
\expandafter\def\csname BookIndex@T0017\endcsname{fanocu@Fano 簇（Fano variety）}
\expandafter\def\csname BookIndex@T0018\endcsname{frauenfelderduliang@Frauenfelder 度量（Frauenfelder metric）}
\expandafter\def\csname BookIndex@T0019\endcsname{frobeniusdaishu@Frobenius 代数（Frobenius algebra）}
\expandafter\def\csname BookIndex@T0020\endcsname{frobeniusliuxing@Frobenius 流形（Frobenius manifold）}
\expandafter\def\csname BookIndex@T0021\endcsname{fukayafanchou@Fukaya 范畴（Fukaya category）}
\expandafter\def\csname BookIndex@T0022\endcsname{gkshoulian@GK 收敛（GK-convergence）}
\expandafter\def\csname BookIndex@T0023\endcsname{grothendieckknudsenliuxing@Grothendieck--Knudsen 流形（Grothendieck–Knudsen manifold）}
\expandafter\def\csname BookIndex@T0024\endcsname{hartmanwintnerdingli@Hartman--Wintner 定理（Hartman–Wintner theorem）}
\expandafter\def\csname BookIndex@T0025\endcsname{heinzjiqiao@Heinz 技巧（Heinz trick）}
\expandafter\def\csname BookIndex@T0026\endcsname{hodgexingsuanzi@Hodge 星算子（Hodge star）}
\expandafter\def\csname BookIndex@T0027\endcsname{hoferduliang@Hofer 度量（Hofer metric）}
\expandafter\def\csname BookIndex@T0028\endcsname{hoferfanshu@Hofer 范数（Hofer norm）}
\expandafter\def\csname BookIndex@T0029\endcsname{kuranishimoxing@Kuranishi 模型（Kuranishi model）}
\expandafter\def\csname BookIndex@T0030\endcsname{kahlerzhui@Kähler 锥（Kähler cone）}
\expandafter\def\csname BookIndex@T0031\endcsname{landauginzburgshi@Landau--Ginzburg 势（Landau–Ginzburg potential）}
\expandafter\def\csname BookIndex@T0032\endcsname{lefschetzshu@Lefschetz 数（Lefschetz number）}
\expandafter\def\csname BookIndex@T0033\endcsname{lefschetzxianweihua@Lefschetz 纤维化（Lefschetz fibration）}
\expandafter\def\csname BookIndex@T0034\endcsname{levicivitalianluo@Levi–Civita 联络（Levi–Civita connection）}
\expandafter\def\csname BookIndex@T0035\endcsname{lipschitzyu@Lipschitz 域（Lipschitz domain）}
\expandafter\def\csname BookIndex@T0036\endcsname{micallefwhitedingli@Micallef--White 定理（Micallef–White theorem）}
\expandafter\def\csname BookIndex@T0037\endcsname{morsewittenshangbianyuansuanzi@Morse--Witten 上边缘算子（Morse–Witten coboundary operator）}
\expandafter\def\csname BookIndex@T0038\endcsname{mumfordfenliang@Mumford 分量（Mumford component）}
\expandafter\def\csname BookIndex@T0039\endcsname{mumfordshang@Mumford 商（Mumford quotient）}
\expandafter\def\csname BookIndex@T0040\endcsname{mobiusbianhuan@Möbius 变换（Möbius transformation）}
\expandafter\def\csname BookIndex@T0041\endcsname{newlandernirenbergdingli@Newlander--Nirenberg 定理（Newlander–Nirenberg theorem）}
\expandafter\def\csname BookIndex@T0042\endcsname{nijenhuiszhangliang@Nijenhuis 张量（Nijenhuis tensor）}
\expandafter\def\csname BookIndex@T0043\endcsname{noetherhuan@Noether 环（Noetherian ring）}
\expandafter\def\csname BookIndex@T0044\endcsname{polyfold@polyfold（polyfold）}
\expandafter\def\csname BookIndex@T0045\endcsname{riemannhurwitzshu@Riemann--Hurwitz 数（Riemann–Hurwitz number）}
\expandafter\def\csname BookIndex@T0046\endcsname{seidelyuan@Seidel 元（Seidel element）}
\expandafter\def\csname BookIndex@T0047\endcsname{seidelbiaoshi@Seidel 表示（Seidel representation）}
\expandafter\def\csname BookIndex@T0048\endcsname{serreduiou@Serre 对偶（Serre duality）}
\expandafter\def\csname BookIndex@T0049\endcsname{severicu@Severi 簇（Severi variety）}
\expandafter\def\csname BookIndex@T0050\endcsname{todajingge@Toda 晶格（Toda lattice）}
\expandafter\def\csname BookIndex@T0051\endcsname{aiermitecongdui@埃尔米特丛对（Hermitian bundle pair）}
\expandafter\def\csname BookIndex@T0052\endcsname{aiermitelianluo@埃尔米特联络（Hermitian connection）}
\expandafter\def\csname BookIndex@T0053\endcsname{banaheliuxing@巴拿赫流形（Banach manifold）}
\expandafter\def\csname BookIndex@T0054\endcsname{bansuibudengshi@伴随不等式（adjunction inequality）}
\expandafter\def\csname BookIndex@T0055\endcsname{bansuigongshi@伴随公式（adjunction formula）}
\expandafter\def\csname BookIndex@T0056\endcsname{bandandiao@半单调（semimonotone）}
\expandafter\def\csname BookIndex@T0057\endcsname{bandanxing@半单性（semisimplicity）}
\expandafter\def\csname BookIndex@T0058\endcsname{banzheng@半正（semipositive）}
\expandafter\def\csname BookIndex@T0059\endcsname{banziyou@半自由（semifree）}
\expandafter\def\csname BookIndex@T0060\endcsname{benyuanziji@本原子集（primitive subset）}
\expandafter\def\csname BookIndex@T0061\endcsname{bianduijiaoxian@边对角线（edge diagonal）}
\expandafter\def\csname BookIndex@T0062\endcsname{bianjiemasiluofuzhibiao@边界马斯洛夫指标（boundary Maslov index）}
\expandafter\def\csname BookIndex@T0063\endcsname{bianjieyingshe@边界映射（boundary map (of flags)）}
\expandafter\def\csname BookIndex@T0064\endcsname{bianlian@边链（edge chain）}
\expandafter\def\csname BookIndex@T0065\endcsname{bianqiuzhiyingshe@边求值映射（edge evaluation map）}
\expandafter\def\csname BookIndex@T0066\endcsname{biaojidian@标记点（marked point）}
\expandafter\def\csname BookIndex@T0067\endcsname{biaojia@标架（framing）}
\expandafter\def\csname BookIndex@T0068\endcsname{bingheziyouji@并合自由积（amalgamated free product）}
\expandafter\def\csname BookIndex@T0069\endcsname{budaicanshudequxian@不带参数的曲线（unparametrized curve）}
\expandafter\def\csname BookIndex@T0070\endcsname{bukeyuequxian@不可约曲线（irreducible curve）}
\expandafter\def\csname BookIndex@T0071\endcsname{canshuzhengze@参数正则（parametrically regular）}
\expandafter\def\csname BookIndex@T0072\endcsname{ceng@层（stratum）}
\expandafter\def\csname BookIndex@T0073\endcsname{cenglunfangfa@层论方法（sheaf-theoretic methods）}
\expandafter\def\csname BookIndex@T0074\endcsname{chaoeweishu@超额维数（excess dimension）}
\expandafter\def\csname BookIndex@T0075\endcsname{chaoliuxing@超流形（supermanifold）}
\expandafter\def\csname BookIndex@T0076\endcsname{chiduweijifen@尺度微积分（scale calculus）}
\expandafter\def\csname BookIndex@T0077\endcsname{chuzifangcheng@除子方程（divisor equation）}
\expandafter\def\csname BookIndex@T0078\endcsname{chuanqianggongshi@穿墙公式（wall-crossing formula）}
\expandafter\def\csname BookIndex@T0079\endcsname{cidiaohehanshu@次调和函数（subharmonic function）}
\expandafter\def\csname BookIndex@T0080\endcsname{congdui@丛对（bundle pair）}
\expandafter\def\csname BookIndex@T0081\endcsname{dabuzuobiaoka@达布坐标卡（Darboux chart）}
\expandafter\def\csname BookIndex@T0082\endcsname{daduijiaoxian@大对角线（fat diagonal）}
\expandafter\def\csname BookIndex@T0083\endcsname{daliangzishangtongdiao@大量子上同调（big quantum cohomology）}
\expandafter\def\csname BookIndex@T0084\endcsname{daibiaohaoshu@带标号树（labelled tree）}
\expandafter\def\csname BookIndex@T0085\endcsname{daibiaojiatonghen@带标架同痕（framed isotopy）}
\expandafter\def\csname BookIndex@T0086\endcsname{daijianbuqiumian@带尖部球面（spiked sphere）}
\expandafter\def\csname BookIndex@T0087\endcsname{daijiediandelimanqumian@带节点的黎曼曲面（nodal Riemann surface）}
\expandafter\def\csname BookIndex@T0088\endcsname{daifujiegou@殆复结构（almost complex structure）}
\expandafter\def\csname BookIndex@T0089\endcsname{daixinliuxing@殆辛流形（almost symplectic manifold）}
\expandafter\def\csname BookIndex@T0090\endcsname{dancebubianliang@单侧不变量（one-sided invariant）}
\expandafter\def\csname BookIndex@T0091\endcsname{danchunquxian@单纯曲线（simple curve）}
\expandafter\def\csname BookIndex@T0092\endcsname{dandiaoxingdingli@单调性定理（monotonicity theorem）}
\expandafter\def\csname BookIndex@T0093\endcsname{danshebanjing@单射半径（injectivity radius）}
\expandafter\def\csname BookIndex@T0094\endcsname{danshedian@单射点（injective point）}
\expandafter\def\csname BookIndex@T0095\endcsname{danzhihua@单值化（monodromy）}
\expandafter\def\csname BookIndex@T0096\endcsname{danzhihuapingyi@单值化平移（monodromy shift）}
\expandafter\def\csname BookIndex@T0097\endcsname{danzhihuashuju@单值化数据（monodromy data）}
\expandafter\def\csname BookIndex@T0098\endcsname{dengwenzuobiao@等温坐标（isothermal coordinates）}
\expandafter\def\csname BookIndex@T0099\endcsname{dengzhoubudengshi@等周不等式（isoperimetric inequality）}
\expandafter\def\csname BookIndex@T0100\endcsname{deliwenti@地理问题（geography problem）}
\expandafter\def\csname BookIndex@T0101\endcsname{dierkeshu@第二可数（second countable）}
\expandafter\def\csname BookIndex@T0102\endcsname{diyikeshu@第一可数（first countable）}
\expandafter\def\csname BookIndex@T0103\endcsname{dianfancong@典范丛（canonical bundle）}
\expandafter\def\csname BookIndex@T0104\endcsname{diaohehanshu@调和函数（harmonic function）}
\expandafter\def\csname BookIndex@T0105\endcsname{diaoheyingshe@调和映射（harmonic map）}
\expandafter\def\csname BookIndex@T0106\endcsname{dingceng@顶层（top stratum）}
\expandafter\def\csname BookIndex@T0107\endcsname{dingxiangpeibian@定向配边（oriented cobordism）}
\expandafter\def\csname BookIndex@T0108\endcsname{duiouhuan@对偶环（dual ring）}
\expandafter\def\csname BookIndex@T0109\endcsname{duiouliangzixishuhuan@对偶量子系数环（dual quantum coefficient ring）}
\expandafter\def\csname BookIndex@T0110\endcsname{duoyuanpan@多圆盘（polydisc）}
\expandafter\def\csname BookIndex@T0111\endcsname{duozhongcidiaohehanshu@多重次调和函数（plurisubharmonic function）}
\expandafter\def\csname BookIndex@T0112\endcsname{duozhongduijiaoxian@多重对角线（multidiagonal）}
\expandafter\def\csname BookIndex@T0113\endcsname{duozhongfugaiquxian@多重覆盖曲线（multiply covered curve）}
\expandafter\def\csname BookIndex@T0114\endcsname{duozhongwending@多重稳定（polystable）}
\expandafter\def\csname BookIndex@T0115\endcsname{erciguji@二次估计（quadratic estimate）}
\expandafter\def\csname BookIndex@T0116\endcsname{erzhongdian@二重点（double point）}
\expandafter\def\csname BookIndex@T0117\endcsname{facong@法丛（normal bundle）}
\expandafter\def\csname BookIndex@T0118\endcsname{faxiangzuobiao@法向坐标（normal coordinate）}
\expandafter\def\csname BookIndex@T0119\endcsname{feidanshedian@非单射点（non-injective point）}
\expandafter\def\csname BookIndex@T0120\endcsname{feijiyadingli@非挤压定理（nonsqueezing theorem）}
\expandafter\def\csname BookIndex@T0121\endcsname{feiqiyidian@非奇异点（nonsingular point）}
\expandafter\def\csname BookIndex@T0122\endcsname{feixinqujian@非辛区间（non-symplectic interval）}
\expandafter\def\csname BookIndex@T0123\endcsname{fenbu@分布（distribution）}
\expandafter\def\csname BookIndex@T0124\endcsname{fengewangluo@分割网络（network of splittings）}
\expandafter\def\csname BookIndex@T0125\endcsname{fenliegongli@分裂公理（splitting axiom）}
\expandafter\def\csname BookIndex@T0126\endcsname{fenshixianxingbianhuan@分式线性变换（fractional linear transformation）}
\expandafter\def\csname BookIndex@T0127\endcsname{fenzhijiechu@分支接触（contact between branches）}
\expandafter\def\csname BookIndex@T0128\endcsname{fuleidehuomusuanzi@弗雷德霍姆算子（Fredholm operator）}
\expandafter\def\csname BookIndex@T0129\endcsname{fuleidehuomuzhibiao@弗雷德霍姆指标（Fredholm index）}
\expandafter\def\csname BookIndex@T0130\endcsname{fuleixieliuxing@弗雷歇流形（Fréchet manifold）}
\expandafter\def\csname BookIndex@T0131\endcsname{fuluoerlianjieguidao@弗洛尔连接轨道（Floer connecting orbit）}
\expandafter\def\csname BookIndex@T0132\endcsname{fuluoertongdiao@弗洛尔同调（Floer homology）}
\expandafter\def\csname BookIndex@T0133\endcsname{fuluoeryantayingshe@弗洛尔延拓映射（Floer continuation map）}
\expandafter\def\csname BookIndex@T0134\endcsname{fubiaojia@复标架（complex frame）}
\expandafter\def\csname BookIndex@T0135\endcsname{fuguixing@复轨形（complex orbifold）}
\expandafter\def\csname BookIndex@T0136\endcsname{ganjingxiangjiao@干净相交（clean intersection）}
\expandafter\def\csname BookIndex@T0137\endcsname{gaojiekuandu@高阶宽度（higher order width）}
\expandafter\def\csname BookIndex@T0138\endcsname{geluomofubi@格罗莫夫闭（Gromov closed）}
\expandafter\def\csname BookIndex@T0139\endcsname{geluomofujinxing@格罗莫夫紧性（Gromov compactness）}
\expandafter\def\csname BookIndex@T0140\endcsname{geluomofukai@格罗莫夫开（Gromov open）}
\expandafter\def\csname BookIndex@T0141\endcsname{geluomofushoulian@格罗莫夫收敛（Gromov convergence）}
\expandafter\def\csname BookIndex@T0142\endcsname{geluomofutapu@格罗莫夫拓扑（Gromov topology）}
\expandafter\def\csname BookIndex@T0143\endcsname{geluomofuweitengbubianliang@格罗莫夫–威滕不变量（Gromov–Witten invariant）}
\expandafter\def\csname BookIndex@T0144\endcsname{geluomofuweitengshi@格罗莫夫--威滕势（Gromov–Witten potential）}
\expandafter\def\csname BookIndex@T0145\endcsname{geluomofuweitengweixunhuan@格罗莫夫--威滕伪循环（Gromov--Witten pseudocycle）}
\expandafter\def\csname BookIndex@T0146\endcsname{genlixiang@根理想（radical ideal）}
\expandafter\def\csname BookIndex@T0147\endcsname{jingbu@颈部（neck）}
\expandafter\def\csname BookIndex@T0148\endcsname{gongxingzhongbiaodu@共形重标度（conformal rescaling）}
\expandafter\def\csname BookIndex@T0149\endcsname{gongxingzuobiao@共形坐标（conformal coordinates）}
\expandafter\def\csname BookIndex@T0150\endcsname{gudingbiaojidian@固定标记点（fixed marked point）}
\expandafter\def\csname BookIndex@T0151\endcsname{guanghuahuasuanzi@光滑化算子（smoothing operator）}
\expandafter\def\csname BookIndex@T0152\endcsname{guangyidehnniuzhuan@广义 Dehn 扭转（generalized Dehn twist）}
\expandafter\def\csname BookIndex@T0153\endcsname{guifandengjia@规范等价（gauge equivalent）}
\expandafter\def\csname BookIndex@T0154\endcsname{guifanhuageluomofuweitengbubianliang@规范化格罗莫夫--威滕不变量（gauged Gromov–Witten invariant）}
\expandafter\def\csname BookIndex@T0155\endcsname{hamidunlianluo@哈密顿联络（Hamiltonian connection）}
\expandafter\def\csname BookIndex@T0156\endcsname{hamidunxianweihua@哈密顿纤维化（Hamiltonian fibration）}
\expandafter\def\csname BookIndex@T0157\endcsname{hamidunraodong@哈密顿扰动（Hamiltonian perturbation）}
\expandafter\def\csname BookIndex@T0158\endcsname{hamiduntonghen@哈密顿同痕（Hamiltonian isotopy）}
\expandafter\def\csname BookIndex@T0159\endcsname{hamidunweifentongpei@哈密顿微分同胚（Hamiltonian diffeomorphism）}
\expandafter\def\csname BookIndex@T0160\endcsname{hamidunxiangliangchang@哈密顿向量场（Hamiltonian vector field）}
\expandafter\def\csname BookIndex@T0161\endcsname{hamidunxinweifentongpei@哈密顿辛微分同胚（Hamiltonian symplectomorphism）}
\expandafter\def\csname BookIndex@T0162\endcsname{xinglieshixian@行列式线（determinant line）}
\expandafter\def\csname BookIndex@T0163\endcsname{haosiduofu@豪斯多夫（Hausdorff）}
\expandafter\def\csname BookIndex@T0164\endcsname{hele@和乐（holonomy）}
\expandafter\def\csname BookIndex@T0165\endcsname{hengjiejiaodian@横截交点（transverse intersection）}
\expandafter\def\csname BookIndex@T0166\endcsname{hengjiexing@横截性（transversality）}
\expandafter\def\csname BookIndex@T0167\endcsname{hengjiezijiaodian@横截自交点（transverse self-intersection）}
\expandafter\def\csname BookIndex@T0168\endcsname{houyishi@后裔势（descendant potential）}
\expandafter\def\csname BookIndex@T0169\endcsname{huanmianliuxing@环面流形（toric manifold）}
\expandafter\def\csname BookIndex@T0170\endcsname{huanyu@环域（annulus）}
\expandafter\def\csname BookIndex@T0171\endcsname{huanyucanshu@环域参数（annulus parameter）}
\expandafter\def\csname BookIndex@T0172\endcsname{hunhebubianliang@混合不变量（mixed invariant）}
\expandafter\def\csname BookIndex@T0173\endcsname{jibenjie@基本解（fundamental solution）}
\expandafter\def\csname BookIndex@T0174\endcsname{jidingli@迹定理（trace theorem）}
\expandafter\def\csname BookIndex@T0175\endcsname{jixiaoqumian@极小曲面（minimal surface）}
\expandafter\def\csname BookIndex@T0176\endcsname{jiyuequxian@既约曲线（reduced curve）}
\expandafter\def\csname BookIndex@T0177\endcsname{jiaquanfanshu@加权范数（weighted norm）}
\expandafter\def\csname BookIndex@T0178\endcsname{jiaquanshu@加权树（weighted tree）}
\expandafter\def\csname BookIndex@T0179\endcsname{jiaquansuoboliefufanshu@加权索伯列夫范数（weighted Sobolev norm）}
\expandafter\def\csname BookIndex@T0180\endcsname{jiaquanyuwendingyingshe@加权预稳定映射（weighted prestable map）}
\expandafter\def\csname BookIndex@T0181\endcsname{jiandianquxian@尖点曲线（cusp-curve）}
\expandafter\def\csname BookIndex@T0182\endcsname{jianpan@尖盘（spiked disc）}
\expandafter\def\csname BookIndex@T0183\endcsname{jiandanceng@简单层（simple stratum）}
\expandafter\def\csname BookIndex@T0184\endcsname{jiandanwendingyingshe@简单稳定映射（simple stable map）}
\expandafter\def\csname BookIndex@T0185\endcsname{jiaobi@交比（cross ratio）}
\expandafter\def\csname BookIndex@T0186\endcsname{jiaodianzhengxing@交点正性（positivity of intersections）}
\expandafter\def\csname BookIndex@T0187\endcsname{jiechuxing@接触型（contact type）}
\expandafter\def\csname BookIndex@T0188\endcsname{jieduanhanshu@截断函数（cutoff function）}
\expandafter\def\csname BookIndex@T0189\endcsname{jiemianlei@截面类（section class）}
\expandafter\def\csname BookIndex@T0190\endcsname{jiemianqulu@截面曲率（sectional curvature）}
\expandafter\def\csname BookIndex@T0191\endcsname{jiedian@节点（nodal point）}
\expandafter\def\csname BookIndex@T0192\endcsname{jinxing@紧性（compactness）}
\expandafter\def\csname BookIndex@T0193\endcsname{jinsijquanchunquxian@近似 $J$-全纯曲线（approximate J-holomorphic curve）}
\expandafter\def\csname BookIndex@T0194\endcsname{jinsidiaohehanshu@近似调和函数（approximately harmonic function）}
\expandafter\def\csname BookIndex@T0195\endcsname{jinsijubujidadian@近似局部极大点（approximate local maximum）}
\expandafter\def\csname BookIndex@T0196\endcsname{jinsiyouni@近似右逆（approximate right inverse）}
\expandafter\def\csname BookIndex@T0197\endcsname{jingqueyouni@精确右逆（exact right inverse）}
\expandafter\def\csname BookIndex@T0198\endcsname{jingxiangzhongcanshuhua@径向重参数化（radial reparametrization）}
\expandafter\def\csname BookIndex@T0199\endcsname{jububiaozhunxing@局部标准形（local normal form）}
\expandafter\def\csname BookIndex@T0200\endcsname{jubuhamidunxianweihua@局部哈密顿纤维化（locally Hamiltonian fibration）}
\expandafter\def\csname BookIndex@T0201\endcsname{jubujiaoshu@局部交数（local intersection number）}
\expandafter\def\csname BookIndex@T0202\endcsname{jubuxinzuoyongliang@局部辛作用量（local symplectic action）}
\expandafter\def\csname BookIndex@T0203\endcsname{jufuluoertongdiao@矩弗洛尔同调（moment Floer homology）}
\expandafter\def\csname BookIndex@T0204\endcsname{juyingshe@矩映射（moment map）}
\expandafter\def\csname BookIndex@T0205\endcsname{juanji@卷积（convolution）}
\expandafter\def\csname BookIndex@T0206\endcsname{jueduilinyushousuohe@绝对邻域收缩核（absolute neighborhood retract）}
\expandafter\def\csname BookIndex@T0207\endcsname{kaileiliuxing@凯勒流形（Kähler manifold）}
\expandafter\def\csname BookIndex@T0208\endcsname{qianrujiechutongdiao@嵌入接触同调（embedded contact homology）}
\expandafter\def\csname BookIndex@T0209\endcsname{kexilimanfangcheng@柯西–黎曼方程（Cauchy–Riemann equation）}
\expandafter\def\csname BookIndex@T0210\endcsname{kexilimansuanzi@柯西–黎曼算子（Cauchy–Riemann operator）}
\expandafter\def\csname BookIndex@T0211\endcsname{kebianbiaojidian@可变标记点（varying marked point）}
\expandafter\def\csname BookIndex@T0212\endcsname{kedulianghua@可度量化（metrizable）}
\expandafter\def\csname BookIndex@T0213\endcsname{kefen@可分（separable）}
\expandafter\def\csname BookIndex@T0214\endcsname{kequqidiandingli@可去奇点定理（removable singularity theorem）}
\expandafter\def\csname BookIndex@T0215\endcsname{keweiyi@可位移（displaceable）}
\expandafter\def\csname BookIndex@T0216\endcsname{keyuequxian@可约曲线（reducible curve）}
\expandafter\def\csname BookIndex@T0217\endcsname{kuxingchengji@裤形乘积（pair-of-pants product）}
\expandafter\def\csname BookIndex@T0218\endcsname{kuandelagelangriziliuxing@宽的拉格朗日子流形（wide Lagrangian）}
\expandafter\def\csname BookIndex@T0219\endcsname{lagelangriziliuxing@拉格朗日子流形（Lagrangian submanifold）}
\expandafter\def\csname BookIndex@T0220\endcsname{lapulasisuanzi@拉普拉斯算子（Laplace operator）}
\expandafter\def\csname BookIndex@T0221\endcsname{limanluohedingli@黎曼--罗赫定理（Riemann–Roch theorem）}
\expandafter\def\csname BookIndex@T0222\endcsname{limanyingshedingli@黎曼映射定理（Riemann mapping theorem）}
\expandafter\def\csname BookIndex@T0223\endcsname{lianluoshi@联络势（connection potential）}
\expandafter\def\csname BookIndex@T0224\endcsname{lianluoxingshi@联络形式（connection form）}
\expandafter\def\csname BookIndex@T0225\endcsname{liantongdanchunxitong@连通单纯系统（connected simple system）}
\expandafter\def\csname BookIndex@T0226\endcsname{liantongdui@连通对（connected pair）}
\expandafter\def\csname BookIndex@T0227\endcsname{liantonghe@连通和（connected sum）}
\expandafter\def\csname BookIndex@T0228\endcsname{liangzibeiji@量子杯积（quantum cup product）}
\expandafter\def\csname BookIndex@T0229\endcsname{liangzishangtongdiao@量子上同调（quantum cohomology）}
\expandafter\def\csname BookIndex@T0230\endcsname{liangzitongdiao@量子同调（quantum homology）}
\expandafter\def\csname BookIndex@T0231\endcsname{liangzixishuhuan@量子系数环（quantum coefficient ring）}
\expandafter\def\csname BookIndex@T0232\endcsname{liangzistanleyreisnerhuan@量子 Stanley--Reisner 环（quantum Stanley–Reisner ring）}
\expandafter\def\csname BookIndex@T0233\endcsname{linjiedian@临界点（critical point）}
\expandafter\def\csname BookIndex@T0234\endcsname{linjiezhi@临界值（critical value）}
\expandafter\def\csname BookIndex@T0235\endcsname{linjinlagelangricaixiang@邻近拉格朗日猜想（nearby Lagrangian conjecture）}
\expandafter\def\csname BookIndex@T0236\endcsname{lingdiandingli@零点定理（Nullstellensatz）}
\expandafter\def\csname BookIndex@T0237\endcsname{liuweierxiangliangchang@刘维尔向量场（Liouville vector field）}
\expandafter\def\csname BookIndex@T0238\endcsname{lujingtishengxingzhi@路径提升性质（path lifting property）}
\expandafter\def\csname BookIndex@T0239\endcsname{luguo@滤过（filtration）}
\expandafter\def\csname BookIndex@T0240\endcsname{luguofuluoertongdiao@滤过弗洛尔同调（filtered Floer homology）}
\expandafter\def\csname BookIndex@T0241\endcsname{masiluofulei@马斯洛夫类（Maslov class）}
\expandafter\def\csname BookIndex@T0242\endcsname{masiluofushu@马斯洛夫数（Maslov number）}
\expandafter\def\csname BookIndex@T0243\endcsname{maopao@冒泡（bubbling）}
\expandafter\def\csname BookIndex@T0244\endcsname{maopaoxianxiang@冒泡现象（bubbling phenomenon）}
\expandafter\def\csname BookIndex@T0245\endcsname{maoji@帽积（cap product）}
\expandafter\def\csname BookIndex@T0246\endcsname{mokongjian@模空间（moduli space）}
\expandafter\def\csname BookIndex@T0247\endcsname{momaopaoshoulian@模冒泡收敛（convergence modulo bubbling）}
\expandafter\def\csname BookIndex@T0248\endcsname{naolu@挠率（torsion）}
\expandafter\def\csname BookIndex@T0249\endcsname{nenglianghengdengshi@能量恒等式（energy identity）}
\expandafter\def\csname BookIndex@T0250\endcsname{nengliangjizhong@能量集中（energy concentration）}
\expandafter\def\csname BookIndex@T0251\endcsname{nengliangmidu@能量密度（energy density）}
\expandafter\def\csname BookIndex@T0252\endcsname{nini@拟逆（pseudoinverse）}
\expandafter\def\csname BookIndex@T0253\endcsname{nitongtai@拟同态（quasimorphism）}
\expandafter\def\csname BookIndex@T0254\endcsname{zhanhe@粘合（gluing）}
\expandafter\def\csname BookIndex@T0255\endcsname{zhanhecanshu@粘合参数（gluing parameter）}
\expandafter\def\csname BookIndex@T0256\endcsname{zhanhedingli@粘合定理（gluing theorem）}
\expandafter\def\csname BookIndex@T0257\endcsname{zhanheyingshe@粘合映射（gluing map）}
\expandafter\def\csname BookIndex@T0258\endcsname{niuqusanciquxian@扭曲三次曲线（twisted cubic）}
\expandafter\def\csname BookIndex@T0259\endcsname{niuqusuanzi@扭曲算子（twisting operator）}
\expandafter\def\csname BookIndex@T0260\endcsname{ouhexingshi@耦合形式（coupling form）}
\expandafter\def\csname BookIndex@T0261\endcsname{pangjialaibudengshi@庞加莱不等式（Poincaré inequality）}
\expandafter\def\csname BookIndex@T0262\endcsname{pengzhangzipingyi@膨胀子平移（dilaton shift）}
\expandafter\def\csname BookIndex@T0263\endcsname{pingxingyidong@平行移动（parallel transport）}
\expandafter\def\csname BookIndex@T0264\endcsname{pinghengxingshi@平衡形式（balanced form）}
\expandafter\def\csname BookIndex@T0265\endcsname{pinghengyingshe@平衡映射（balanced map）}
\expandafter\def\csname BookIndex@T0266\endcsname{pingjunzhibudengshi@平均值不等式（mean value inequality）}
\expandafter\def\csname BookIndex@T0267\endcsname{pingjunzhixingzhi@平均值性质（mean value property）}
\expandafter\def\csname BookIndex@T0268\endcsname{pushimokongjian@普适模空间（universal moduli space）}
\expandafter\def\csname BookIndex@T0269\endcsname{pushiqiuzhiyingshe@普适求值映射（universal evaluation map）}
\expandafter\def\csname BookIndex@T0270\endcsname{pushinovikovhuan@普适 Novikov 环（universal Novikov ring）}
\expandafter\def\csname BookIndex@T0271\endcsname{pububianliang@谱不变量（spectral invariant）}
\expandafter\def\csname BookIndex@T0272\endcsname{pufanshu@谱范数（spectral norm）}
\expandafter\def\csname BookIndex@T0273\endcsname{puliu@谱流（spectral flow）}
\expandafter\def\csname BookIndex@T0274\endcsname{qi@旗（flag (of a tree)）}
\expandafter\def\csname BookIndex@T0275\endcsname{qipao@气泡（bubble）}
\expandafter\def\csname BookIndex@T0276\endcsname{qipaonengliang@气泡能量（bubble energy）}
\expandafter\def\csname BookIndex@T0277\endcsname{qiadangyingshe@恰当映射（proper map）}
\expandafter\def\csname BookIndex@T0278\endcsname{xianweihe@纤维和（fiber sum）}
\expandafter\def\csname BookIndex@T0279\endcsname{xianweihua@纤维化（fibration）}
\expandafter\def\csname BookIndex@T0280\endcsname{qiang@腔（chamber）}
\expandafter\def\csname BookIndex@T0281\endcsname{qianghengjie@强横截（strongly transverse）}
\expandafter\def\csname BookIndex@T0282\endcsname{qiezhui@切锥（tangent cone）}
\expandafter\def\csname BookIndex@T0283\endcsname{qiuzhiyingshe@求值映射（evaluation map）}
\expandafter\def\csname BookIndex@T0284\endcsname{qiumiandandiao@球面单调（spherically monotone）}
\expandafter\def\csname BookIndex@T0285\endcsname{qiumianlei@球面类（spherical class）}
\expandafter\def\csname BookIndex@T0286\endcsname{quancede@全测地（totally geodesic）}
\expandafter\def\csname BookIndex@T0287\endcsname{quanchunfenqifugaiyingshe@全纯分歧覆盖映射（holomorphic branched covering）}
\expandafter\def\csname BookIndex@T0288\endcsname{quanchunjiegou@全纯结构（holomorphic structure）}
\expandafter\def\csname BookIndex@T0289\endcsname{quanjujiaoxingshi@全局角形式（global angular form）}
\expandafter\def\csname BookIndex@T0290\endcsname{quanshiziliuxing@全实子流形（totally real submanifold）}
\expandafter\def\csname BookIndex@T0291\endcsname{quanhanshu@权函数（weight function）}
\expandafter\def\csname BookIndex@T0292\endcsname{ruodaibiaoyuan@弱代表元（weak representative）}
\expandafter\def\csname BookIndex@T0293\endcsname{ruodaoshu@弱导数（weak derivative）}
\expandafter\def\csname BookIndex@T0294\endcsname{ruohengjie@弱横截（weakly transverse）}
\expandafter\def\csname BookIndex@T0295\endcsname{ruojie@弱解（weak solution）}
\expandafter\def\csname BookIndex@T0296\endcsname{sadedingli@萨德定理（Sard theorem）}
\expandafter\def\csname BookIndex@T0297\endcsname{senlin@森林（forest）}
\expandafter\def\csname BookIndex@T0298\endcsname{shangjingtanyingshe@上惊叹映射（upper shriek map）}
\expandafter\def\csname BookIndex@T0299\endcsname{shangtongdiaochanglun@上同调场论（Cohomological Field Theory）}
\expandafter\def\csname BookIndex@T0300\endcsname{shengyuji@剩余集（residual set）}
\expandafter\def\csname BookIndex@T0301\endcsname{shihanshu@势函数（potential function）}
\expandafter\def\csname BookIndex@T0302\endcsname{shuzhiyouxiaonef@数值有效（NEF）（numerically effective (NEF)）}
\expandafter\def\csname BookIndex@T0303\endcsname{shububianliang@树不变量（invariant for trees）}
\expandafter\def\csname BookIndex@T0304\endcsname{shutonggou@树同构（tree isomorphism）}
\expandafter\def\csname BookIndex@T0305\endcsname{shutongtai@树同态（tree homomorphism）}
\expandafter\def\csname BookIndex@T0306\endcsname{shuzhidaifujiegou@竖直殆复结构（vertical almost complex structure）}
\expandafter\def\csname BookIndex@T0307\endcsname{shuzhinengliang@竖直能量（vertical energy）}
\expandafter\def\csname BookIndex@T0308\endcsname{shuzhiqiecong@竖直切丛（vertical tangent bundle）}
\expandafter\def\csname BookIndex@T0309\endcsname{shuzhiqiumianlei@竖直球面类（vertical spherical class）}
\expandafter\def\csname BookIndex@T0310\endcsname{shuzhiweifen@竖直微分（vertical differential）}
\expandafter\def\csname BookIndex@T0311\endcsname{shuzhizhengze@竖直正则（vertically regular）}
\expandafter\def\csname BookIndex@T0312\endcsname{shuzhilevicivitalianluo@竖直 Levi--Civita 联络（vertical Levi–Civita connection）}
\expandafter\def\csname BookIndex@T0313\endcsname{shuangbubianfinslerduliang@双不变 Finsler 度量（bi-invariant Finsler metric）}
\expandafter\def\csname BookIndex@T0314\endcsname{shuipingfenbu@水平分布（horizontal distribution）}
\expandafter\def\csname BookIndex@T0315\endcsname{shuipingjiemian@水平截面（horizontal section）}
\expandafter\def\csname BookIndex@T0316\endcsname{shuipinglei@水平类（horizontal class）}
\expandafter\def\csname BookIndex@T0317\endcsname{shuipingtisheng@水平提升（horizontal lift）}
\expandafter\def\csname BookIndex@T0318\endcsname{suoboliefuqianrudingli@索伯列夫嵌入定理（Sobolev embedding theorem）}
\expandafter\def\csname BookIndex@T0319\endcsname{suoboliefukongjian@索伯列夫空间（Sobolev space）}
\expandafter\def\csname BookIndex@T0320\endcsname{teshudian@特殊点（special point）}
\expandafter\def\csname BookIndex@T0321\endcsname{teshudingdian@特殊顶点（special vertex）}
\expandafter\def\csname BookIndex@T0322\endcsname{tongdiaobenzhilinjiezhi@同调本质临界值（homologically essential critical value）}
\expandafter\def\csname BookIndex@T0323\endcsname{tongdiaolahui@同调拉回（homology pullback）}
\expandafter\def\csname BookIndex@T0324\endcsname{tuxiangweixunhuan@图像伪循环（graph pseudocycle）}
\expandafter\def\csname BookIndex@T0325\endcsname{tuxiangyingshe@图像映射（graph map）}
\expandafter\def\csname BookIndex@T0326\endcsname{tuichu@推出（pushout）}
\expandafter\def\csname BookIndex@T0327\endcsname{tuoyuanguji@椭圆估计（elliptic estimates）}
\expandafter\def\csname BookIndex@T0328\endcsname{tuoyuanqumian@椭圆曲面（elliptic surface）}
\expandafter\def\csname BookIndex@T0329\endcsname{tuoyuanzhengzexing@椭圆正则性（elliptic regularity）}
\expandafter\def\csname BookIndex@T0330\endcsname{tuoyuanziju@椭圆自举（elliptic bootstrapping）}
\expandafter\def\csname BookIndex@T0331\endcsname{waibeiji@外杯积（external cup product）}
\expandafter\def\csname BookIndex@T0332\endcsname{weiyifenjiezhenghuan@唯一分解整环（unique factorization domain）}
\expandafter\def\csname BookIndex@T0333\endcsname{weiyiyanta@唯一延拓（unique continuation）}
\expandafter\def\csname BookIndex@T0334\endcsname{weiquanchunquxian@伪全纯曲线（pseudoholomorphic curve）}
\expandafter\def\csname BookIndex@T0335\endcsname{weiquanchuntuxiang@伪全纯图像（pseudoholomorphic graph）}
\expandafter\def\csname BookIndex@T0336\endcsname{weitongdiao@伪同调（pseudohomology）}
\expandafter\def\csname BookIndex@T0337\endcsname{weixunhuan@伪循环（pseudocycle）}
\expandafter\def\csname BookIndex@T0338\endcsname{wendingfanchou@稳定范畴（stable category）}
\expandafter\def\csname BookIndex@T0339\endcsname{wendingquxian@稳定曲线（stable curve）}
\expandafter\def\csname BookIndex@T0340\endcsname{wendingyingshe@稳定映射（stable map）}
\expandafter\def\csname BookIndex@T0341\endcsname{wendingyingshetapu@稳定映射拓扑（stable map topology）}
\expandafter\def\csname BookIndex@T0342\endcsname{wendingziqun@稳定子群（isotropy subgroup）}
\expandafter\def\csname BookIndex@T0343\endcsname{wuqiongxiaodaifujiegou@无穷小殆复结构（infinitesimal almost complex structure）}
\expandafter\def\csname BookIndex@T0344\endcsname{wuqiongyuanchutu@无穷远处凸（convex at infinity）}
\expandafter\def\csname BookIndex@T0345\endcsname{xianfangcheng@弦方程（string equation）}
\expandafter\def\csname BookIndex@T0346\endcsname{xianxinghuasuanzi@线性化算子（linearized operator）}
\expandafter\def\csname BookIndex@T0347\endcsname{xiangduixinzuoyongliang@相对辛作用量（relative symplectic action）}
\expandafter\def\csname BookIndex@T0348\endcsname{xiangrong@相容（compatible）}
\expandafter\def\csname BookIndex@T0349\endcsname{xiaoliangzichengfa@小量子乘法（small quantum multiplication）}
\expandafter\def\csname BookIndex@T0350\endcsname{xiebansui@斜伴随（skew-adjoint）}
\expandafter\def\csname BookIndex@T0351\endcsname{xinbu@辛补（symplectic complement）}
\expandafter\def\csname BookIndex@T0352\endcsname{xinchanglun@辛场论（symplectic field theory）}
\expandafter\def\csname BookIndex@T0353\endcsname{xinchuixia@辛吹下（symplectic blowdown）}
\expandafter\def\csname BookIndex@T0354\endcsname{xindanyouli@辛单有理（symplectically uniruled）}
\expandafter\def\csname BookIndex@T0355\endcsname{xinfeiqiumian@辛非球面（symplectically aspherical）}
\expandafter\def\csname BookIndex@T0356\endcsname{xinwoxuanfangcheng@辛涡旋方程（symplectic vortex equation）}
\expandafter\def\csname BookIndex@T0357\endcsname{xinkuandu@辛宽度（symplectic width）}
\expandafter\def\csname BookIndex@T0358\endcsname{xinliuxing@辛流形（symplectic manifold）}
\expandafter\def\csname BookIndex@T0359\endcsname{xinpengzhang@辛膨胀（symplectic inflation）}
\expandafter\def\csname BookIndex@T0360\endcsname{xinqiege@辛切割（symplectic cutting）}
\expandafter\def\csname BookIndex@T0361\endcsname{xintonghen@辛同痕（symplectically isotopic）}
\expandafter\def\csname BookIndex@T0362\endcsname{xinweifentongpeiqun@辛微分同胚群（symplectomorphism group）}
\expandafter\def\csname BookIndex@T0363\endcsname{xinxingshi@辛形式（symplectic form）}
\expandafter\def\csname BookIndex@T0364\endcsname{xinxingshixingbian@辛形式形变（symplectic deformation）}
\expandafter\def\csname BookIndex@T0365\endcsname{xinyingsheleiqun@辛映射类群（symplectic mapping class group）}
\expandafter\def\csname BookIndex@T0366\endcsname{xinzhuangtian@辛装填（symplectic packing）}
\expandafter\def\csname BookIndex@T0367\endcsname{xinzhui@辛锥（symplectic cone）}
\expandafter\def\csname BookIndex@T0368\endcsname{xingshibansuisuanzi@形式伴随算子（formal adjoint operator）}
\expandafter\def\csname BookIndex@T0369\endcsname{xujibenxunhuan@虚基本循环（virtual fundamental cycle）}
\expandafter\def\csname BookIndex@T0370\endcsname{xumoxunhuan@虚模循环（virtual moduli cycle）}
\expandafter\def\csname BookIndex@T0371\endcsname{xuliejin@序列紧（sequentially compact）}
\expandafter\def\csname BookIndex@T0372\endcsname{xunfu@驯服（tame）}
\expandafter\def\csname BookIndex@T0373\endcsname{yapingquxian@压平曲线（flattened curve）}
\expandafter\def\csname BookIndex@T0374\endcsname{yanxianweijifen@沿纤维积分（fiber integration）}
\expandafter\def\csname BookIndex@T0375\endcsname{yizhihua@一致化（uniformization）}
\expandafter\def\csname BookIndex@T0376\endcsname{yilaidingyiyudiandedaifujiegou@依赖定义域点的殆复结构（domain-dependent almost complex structure）}
\expandafter\def\csname BookIndex@T0377\endcsname{yiwangyingshe@遗忘映射（forgetful map）}
\expandafter\def\csname BookIndex@T0378\endcsname{yinlihouyi@引力后裔（gravitational descendant）}
\expandafter\def\csname BookIndex@T0379\endcsname{yinlizuxian@引力祖先（gravitational ancestor）}
\expandafter\def\csname BookIndex@T0380\endcsname{yinhanshudingli@隐函数定理（implicit function theorem）}
\expandafter\def\csname BookIndex@T0381\endcsname{yingsheyuan@映射圆（mapping torus）}
\expandafter\def\csname BookIndex@T0382\endcsname{youlingfenzhi@幽灵分支（ghost component）}
\expandafter\def\csname BookIndex@T0383\endcsname{youlingmianpian@幽灵面片（ghost facet）}
\expandafter\def\csname BookIndex@T0384\endcsname{youlingqipao@幽灵气泡（ghost bubble）}
\expandafter\def\csname BookIndex@T0385\endcsname{youliliantong@有理连通（rationally connected）}
\expandafter\def\csname BookIndex@T0386\endcsname{youliquxian@有理曲线（rational curve）}
\expandafter\def\csname BookIndex@T0387\endcsname{youlizhengguiquxian@有理正规曲线（rational normal curve）}
\expandafter\def\csname BookIndex@T0388\endcsname{youxianweiyuehua@有限维约化（finite dimensional reduction）}
\expandafter\def\csname BookIndex@T0389\endcsname{youxiaochuzi@有效除子（effective divisor）}
\expandafter\def\csname BookIndex@T0390\endcsname{youxiaozhui@有效锥（effective cone）}
\expandafter\def\csname BookIndex@T0391\endcsname{youpingfanhua@酉平凡化（unitary trivialization）}
\expandafter\def\csname BookIndex@T0392\endcsname{youni@右逆（right inverse）}
\expandafter\def\csname BookIndex@T0393\endcsname{yuqiaxingshi@余恰形式（coexact form）}
\expandafter\def\csname BookIndex@T0394\endcsname{yuzhanhe@预粘合（pregluing）}
\expandafter\def\csname BookIndex@T0395\endcsname{yuzhanhequxian@预粘合曲线（preglued curve）}
\expandafter\def\csname BookIndex@T0396\endcsname{yuwendingquxian@预稳定曲线（prestable curve）}
\expandafter\def\csname BookIndex@T0397\endcsname{zaimouchudanshe@在某处单射（somewhere injective）}
\expandafter\def\csname BookIndex@T0398\endcsname{zhaidelagelangriziliuxing@窄的拉格朗日子流形（narrow Lagrangian）}
\expandafter\def\csname BookIndex@T0399\endcsname{zhangaicong@障碍丛（obstruction bundle）}
\expandafter\def\csname BookIndex@T0400\endcsname{zhengguizuobiao@正规坐标（normal coordinates）}
\expandafter\def\csname BookIndex@T0401\endcsname{zhenghelagelangriziliuxing@正合拉格朗日子流形（exact Lagrangian submanifold）}
\expandafter\def\csname BookIndex@T0402\endcsname{zhengheraodong@正合扰动（exact perturbation）}
\expandafter\def\csname BookIndex@T0403\endcsname{zhenghexingshi@正合形式（exact form）}
\expandafter\def\csname BookIndex@T0404\endcsname{zhengmianjixingshi@正面积形式（positive area form）}
\expandafter\def\csname BookIndex@T0405\endcsname{zhengzedian@正则点（regular point）}
\expandafter\def\csname BookIndex@T0406\endcsname{zhengzedui@正则对（regular pair）}
\expandafter\def\csname BookIndex@T0407\endcsname{zhengzetonglun@正则同伦（regular homotopy）}
\expandafter\def\csname BookIndex@T0408\endcsname{zhengzexing@正则性（regularity）}
\expandafter\def\csname BookIndex@T0409\endcsname{zhengzezhi@正则值（regular value）}
\expandafter\def\csname BookIndex@T0410\endcsname{zhibiaojifenge@指标集分割（splitting of an index set）}
\expandafter\def\csname BookIndex@T0411\endcsname{zhongbiaoduhoudeyingshelie@重标度后的映射列（renormalized sequence）}
\expandafter\def\csname BookIndex@T0412\endcsname{zhongbiaoduxulie@重标度序列（rescaling sequence）}
\expandafter\def\csname BookIndex@T0413\endcsname{zhongcanshuhuaqun@重参数化群（reparametrization group）}
\expandafter\def\csname BookIndex@T0414\endcsname{zhongyanxiancong@重言线丛（tautological line bundle）}
\expandafter\def\csname BookIndex@T0415\endcsname{zhudianyueshu@逐点约束（pointwise constraint）}
\expandafter\def\csname BookIndex@T0416\endcsname{zhujinyizhishoulian@逐紧一致收敛（uniform convergence on compact subsets (u.c.s.)）}
\expandafter\def\csname BookIndex@T0417\endcsname{zishu@子树（subtree）}
\expandafter\def\csname BookIndex@T0418\endcsname{zuixiaochenshu@最小陈数（minimal Chern number）}
\expandafter\def\csname BookIndex@T0419\endcsname{zuoyongpu@作用谱（action spectrum）}
\makeatother

\let\BookOriginalIndex\index
\renewcommand{\index}[1]{\ifBookIndexing\BookOriginalIndex{#1}\fi}
